\setchapter{7}{Nonparametric and Semiparametric Estimators}
\ChapterHeading

A \Term{parametric estimator} is used to estiamte \KeyIdea{finite-dimensional parameters} of a model where
these parameter values uniquely identify the conditional probability of the dependent variables. 
A \Term{nonparametric estimator} is use to estimate features of a model \Emph{without making any assumptions} on the population distribution.
It is often used to \Emph{estimate functions}, such as the probability density function or a conditional mean without assumming strong conditions
such as normality. These are denoted as \KeyIdea{infinite dimensional parameters}.
A \Term{semiparametric estimator} is one sued to estimate \KeyIdea{finite-dimensional parameters} but these do not uniquely identify the conditional probabilty distribution of the dependent variables.\\

The same estimator can be seen as parametric, nonparametric, and semiparametric depending on the context. Consider the linear regression model 
\begin{equation*}
    y = x'\beta_0+u
\end{equation*}
and the Ordinary Least Squares Estimator (OLS).
\begin{itemize}
    \item OLS is \Emph{parametric} if we assume $u|x\sim N(0,\sigma^2)$ and obtain the estimator via ML.
    \item OLS is \Emph{semiparametric} if we assume $E[u|x] = 0$ thus $E[y|x] = x'\beta_0$. we do not assume normalty of the error term, but assume a specific form for the mean of $y$ conditional on $x$.
    \item OLS is \Emph{nonparametric} if our objective is to estiamte the best linear predictor coefficients. That is, we want to estiamte the value of $\beta$ that minimizes $E[(y-x'\beta)^2]$.
    \Info{We don't assume anything. Just finding the function to minimize $E[(y-x'\beta)^2]$ but no assumption on the distribution of the data.}
\end{itemize}

\section{Nonparametric estimation}

Suppose $\{y_i\}_{i=1}^{n}$ are iid where $y_i$ is a random variable with a continuous distribution.
Our objective is to \Emph{estimate the density function of $y\to f(y)$}.

\subsection{Crude estimator}

A \Term{crude estimator} of the density of $y$ is the histogram.
The shape of the histogram depends on the \Emph{bin width}, on the \Emph{number of bins}, and on the chosen \Emph{origin}.
Fixing these parameters, the density of $y$ at $y=y_0$ given by the histogram is
\begin{equation*}
    \hat{f}(y_0) = \frac{1}{n}\frac{\text{number of observations in the same bin as $y_0$}}{\text{width of the bin containg $y_0$}}
\end{equation*}
If we define $h$ as the bin width and let $y^*$ denote the midpoint of the interval containing $y_0$, the histogram estiamte is given by 
\begin{equation*}
    \hat{f}(y_0) = \frac{1}{nh}\sum_{j=1}^{n}1(-\frac{1}{2}\leq \frac{y_j-y^*}{h}<\frac{1}{2})
\end{equation*}
where $1(A) = \begin{cases}
    1 & \text{if A is true}\\
    0 & \text{if A is not true}
\end{cases} \to $ Bernoulli distribution. Note the estiamtor is \Emph{not continuous at $y_0$}.

\subsection{Naive density estimator}

From the definition of a probability density, we have
\begin{align*}
    f(y_0) & = \lim_{h\to 0}\frac{F(y_0+\frac{h}{2}) - F(y_0-\frac{h}{2})}{h}\\
    & = \lim_{h\to 0}\frac{1}{h} P(y_0-\frac{h}{2}<y<y_0+\frac{h}{2})\\
    & = \lim_{h\to 0}\frac{1}{h} E\left[1\left(y_0-\frac{h}{2}<y<y_0+\frac{h}{2}\right)\to\text{ Bernoulli}\right]\\
    & = 1\times P\left(y_0-\frac{h}{2}\leq y<y_0+\frac{h}{2}\right) + 0\times \left(1-P\left(y_0-\frac{h}{2}<y<y_0+\frac{h}{2}\right)\right)\\
    & = P\left(y_0-\frac{h}{2}\leq y<y_0+\frac{h}{2}\right)
\end{align*}
\Info{Then use the averages to find the expected values.}
Hence the estimator of $f(y_0)$ given by the \Emph{analogy principle} is the \Term{naive density estimator}
\begin{align*}
    \hat{f}(y_0) & = \frac{1}{nh}\sum_{j=1}^{n} 1\left(y_0-\frac{h}{2}\leq y<y_0+\frac{h}{2}\right)\\
    & = \frac{1}{nh}\sum_{j=1}^{n} 1\left(-\frac{1}{2}\leq \frac{y_j-y_0}{h}<y_0+\frac{1}{2}\right)
\end{align*}
For the estimator to be \Emph{consistent and ensures convergence} in iid samples, we need \KeyIdea{$n\to\infty$, $nh\to\infty$ and $h\to 0$}.
\begin{proof}
    $h\to 0$
    \begin{align*}
        E{\hat{f}(y_0)} & = E\left[\frac{1}{nh}\sum_{j=1}^{n} 1\left(y_0-\frac{h}{2}\leq y_j<y_0+\frac{h}{2}\right)\right]\\
        & = \frac{1}{nh}\sum_{j=1}^{n}E\left[1\left(y_0-\frac{h}{2}\leq y<y_0+\frac{h}{2}\right)\right] \text{($y_j\to y$ because iid)}\\
        & = \frac{n}{nh} \underbrace{E\left[1\left(y_0-\frac{h}{2}\leq y<y_0+\frac{h}{2}\right)\right]}_{f(y_0)\text{ definition}}\\
        & = \frac{1}{h}\left[F\left(y_0-\frac{h}{2}\right)-F(y_0-\frac{h}{2})\right]
    \end{align*}
    So, from $f(y_0)$ definition,
    \begin{equation*}
        \frac{1}{h}\left[F\left(y_0-\frac{h}{2}\right)-F(y_0-\frac{h}{2})\right] \to f(y_0) \;\text{ if } h\to 0
    \end{equation*}
\end{proof}
\begin{proof}
    $nh\to\infty$
    \begin{align*}
        Var(\hat{f}(y_0)) & = Var\left[\frac{1}{nh}\sum_{j=1}^{n} 1\left(y_0-\frac{h}{2}\leq y_j<y_0+\frac{h}{2}\right)\right]\\
        & = \frac{1}{n^2h^2}\sum_{j=1}^{n}Var\left[1\left(y_0-\frac{h}{2}\leq y_j<y_0+\frac{h}{2}\right)\right] \text{(variance of sum = sum of variance)}\\
        & = \frac{n}{n^2h^2}\underbrace{Var\left[1\left(y_0-\frac{h}{2}\leq y<y_0+\frac{h}{2}\right)\right]}_{\text{Var[Bernoulli] = $\theta(1-\theta)$}}\\
        & = \frac{1}{nh^2} P\left(y_0-\frac{h}{2}\leq y<y_0+\frac{h}{2}\right) \times \underbrace{\left[1 - P\left(y_0-\frac{h}{2}\leq y<y_0+\frac{h}{2}\right)\right]}_{\leq 1}
    \end{align*}
    and because $0\leq Var(\hat{f}(y_0))$
    \begin{equation*}
        0\leq Var(\hat{f}(y_0))\leq \frac{1}{nh}\underbrace{\frac{P\left(y_0-\frac{h}{2}\leq y<y_0+\frac{h}{2}\right)}{h}}_{f(y_0)\to\text{constant if }h\to 0} 
    \end{equation*}
    Thus we conclude $\frac{1}{nh}\to 0 \Rightarrow nh\to\infty$
\end{proof}
Consequently, $\lim\hat{f}(y_0) \xrightarrow{p} f(y_0)$. 
\Info{If $\lim E[\hat{\theta}] = \theta_0$ and $\lim Var[\hat{\theta}]=0$ then $\lim \hat{\theta}\xrightarrow{p}\theta_0$}.
Like the histogram estimator, the naive estimator is \Emph{not continuous at $y_0$}.

\subsection{Kernel estimator}

A continuous and smooth estimator can be obtained by replacing the function $K[z] = 1\left(-\frac{1}{2}\leq z < \frac{1}{2}\right)$ by another \Emph{kernel function}.

\begin{definition}
    A \textbf{Kernel function} is any function $K(\cdot)$ such that 
    \begin{equation*}
        \int_{\bbR}K(u)\, du = 1
    \end{equation*}
    (Rosenblatt, 1956)
\end{definition}

We do not require $K(\cdot)\geq 0$ although this assumption ensures that the estimator is non-negative.
Many kernels are available but a popular choice is to use \KeyIdea{$K[z] = \phi(z)$}, the density function of the standard normal distribution.
In general, for a bandwidth $h>0$, the \Emph{kernel density function} is given by 
\begin{equation*}
    \hat{f}(y_0) = \frac{1}{nh}\sum_{j=1}^{n}K\left(\frac{y_i-y_0}{h}\right)
\end{equation*}
The choice of the kernel \Emph{affects the bias} $E[\hat{f}(y_0)] - f(y_0)$.
\Emph{Consistency} requires \KeyIdea{$n\to\infty$, $nh\to\infty$, and $h\to 0$}.\\

This can be generalized to d-dimensional data using the multivaraite kernel function $K_d(\cdot)$ such that $\int_{\bbR^d}K_d(u)\,du=1$.
To ensure \Emph{consistency}, we need to assume \KeyIdea{$n\to\infty$, $h\to 0$, and $nh^d\to\infty$}.
In this case the variance of the estimator goes to zero as $nh^d\to\infty$.
Notice the \Emph{curse of dimensionality}, the larger the value of $d$, the slower the rate of convergence.
\Info{The norm $||\hat{\beta}-\beta_0|| \leq \frac{c}{N^{1/2}}\to$ goes to zero much slower than parametric.}
Under additional conditions it is also possible to ensure \Emph{asymptotically normaity} but the rate of convergence is slower than rate $n^{-1/2}$ and the estimator is \Emph{biased}.
The rate of convergence depends on $K_d(\cdot)$ but other authors say this choice is relatively unimportant.
The \Emph{bandwidth of $h$} affects the \Emph{bias and variance} of the estimator, as $h$ decreases, so does the bias; however, the variance increases.\\

The choice of \Emph{bandwidth} is critical for the result.
\Emph{Too large} means \KeyIdea{over-smoothing} and elimiates important features of the function being estimated.
\Emph{Too small} means \KeyIdea{under-smoothing} and the important features of the function become obscured by noise.
Choices of implementation rules include
\begin{itemize}
    \item Automatic rules
    \item Graphically analyzing the effects of varying the bandwidth
    \item Minimizing selection criterion
    \item Silverman's rule of thumb $\to$ popular choice for the Gaussian kernel 
    \begin{equation*}
        h = \min\left\{\sigma,\frac{IQR}{1,34}\right\}\frac{0.9}{n^{1/5}}
    \end{equation*}
\end{itemize}

\subsection{Kernel regression : discrete $x_i$}

The kernel density estimator can be used for \Term{kernel regression}.
It is particularly useful for exploratory analysis of the data as it may be covenient not to assume a functional form for the relation between two variables.\\

Let $\{(y_i,x_i)\}_{i=1}^{n}$ be iid where $y_i$ and $x_i$ are random variables with continuous distributions and suppose we want to estimate $E[y_i|x_i] = m(x_i)$.
We can write this as 
\begin{equation*}
    y_i = m(x_i) + \epsilon_i
\end{equation*}
where $E[\epsilon_i|x_i] = 0$.
\begin{proof}
    $E[\epsilon_i|x_i] = 0\to y_i = m(x_i) + \epsilon_i$
    \begin{align*}
        \epsilon_i & = y_i - E[y_i|x_i]\\
        E[\epsilon_i|x_i] & = E[y_i - E[y_i|x_i]|x_i]\\
        & = E[y_i|x_i] - E[E[y_i|x_i]|x_i]\\
        & = E[y_i|x_i] - E[y_i|x_i] = 0\\
        \Rightarrow y_i & = \underbrace{E[y_i|x_i]}_{m_i} + \epsilon_i 
    \end{align*}
\end{proof}

If we are interested in estimating $m(x_j)$ at a particular value of $x$ and if \Emph{multiple observations of $y$} were available for that particular value of $x$,
we could use
\begin{align*}
    \hat{m}(x_0) & = \frac{\sum_{j=1}^{n}1(x_j=x_0)y_j}{\sum_{j=1}^{n}1(x_j=x_0)} = \sum_{j=1}^{n}\frac{1(x_j=x_0)}{\sum_{j=1}^{n}1(x_j=x_0)}\\
    & = m(x_0) + \frac{\sum_{j=1}^{n}1(x_j=x_0)\epsilon_j}{\sum_{j=1}^{n}1(x_j=x_0)}
\end{align*}
which under mild assumptions is \Emph{consistent} if \KeyIdea{$n\to\infty$} provided that $P(x_j=x_0)>0$. 
\Info{Only assuming we have many observations at $x_0$.}
Consequently, this estimator is \Emph{not valid} if $x_j$ is a continuous random variable.
\Info{If data is continous, $P(x=x_0)=0$, integral at a point has no mass $\to 0 \to$ not valid.}\\

\subsection{Kernel regression : continous $x_i$}

For \Emph{continuous random variables} if $m(\cdot)$ is sufficiently smooth, we can estimate $m(x_0)$ using observations in a \KeyIdea{neighborhood of $x_0$}.
A \Term{smoothing estimator} of $m(x_0)$ can be written as 
\begin{equation*}
    \hat{m}(x_0) = \sum_{j=1}^{n}\omega_h(x_0-x_j)y_j
\end{equation*}
where the weights $\omega_h(x_0-x_j)$ \Emph{decreases} with the distance between between $x_j$ and $x_0$.
Frequently, the weights are defined as 
\begin{equation*}
    \omega_h(x_0-x_j) = \frac{K(\frac{x_0-x_j}{h})}{\sum_{k=1}^{n}K(\frac{x_0-x_j}{h})}
\end{equation*}
This is the \Term{Nadaraya-Watson estimator} of $m(x_0)$.
Under the regularity conditions, $\hat{m}(x_0)$ is \Emph{consistent} when \KeyIdea{$\lim_{n\to\infty}h=0$}, but convergence is slow.\\

The properties of the estimator depend on the choice of $K(\cdot)$ and $h$. 
The choice of $K(\cdot)$ is not that import, often just use $\phi(\cdot)$.
But the \Emph{choice for $h$ is critical}.
If $h$ is too small, the estimates become too irregular. If $h$ is too large, the estimates appraoch $\bar{y}$.
A popular method to choose $h$ is \KeyIdea{leave-one-out cross validation}.
\begin{equation*}
    h_{cv} = \text{arg }\min_{h}\frac{1}{n}\sum_{i=1}^{n}\delta(x_i)\left[y_i - \sum_{j\neq i}\omega_h(x_i-x_0)y_i\right]^2
\end{equation*}
where $\delta(x)$ is a \Emph{trimming function} to reduce boundary effects.
\Info{$\delta(x)$ trimming to leave out the tails, 0 in the tail $\to$ indicator function.}
This can be generalized for a vector of regressors. 
However, for the multivariate case, not only do we have \Emph{curse of dimensionality} but also lose the ability to display results graphically.

\section{Semiparametric estimation}

Most non-ML estimators are semiparametric in the sense that they do not require the full specification of the conditional probability distribution of the dependent variable.

\subsection{Partially linear model}

Suppose $E[y_i|z,x] = z_i\beta + g(x_i)$ where \Emph{$g(\cdot)$ is left unspecified}, and 
\begin{equation*}
    y_i = z_i\beta + g(x_i) + \epsilon_i,\quad i=1,\ldots,n
\end{equation*}
\Term{Robinson (1988) and Speckman (1988)} have shown that under certain conditions it is possible to obtain $\sqrt{n}$-consistent estimators for $\beta$.
Notice we can write 
\begin{align*}
    E[y_i|x] & = E[z_i|x]\beta + g(x_i)\\
    y_i-E[y_i|x] & = [z_i-E[z_i|x]]\beta + \epsilon_i
\end{align*}
\begin{proof}
    We begin with initial condition: $E[\epsilon_i|z_i,x_i] = 0$
    \begin{align*}
        E[y_i|x_i] & = E[z_i\beta+g(x_i)+\epsilon_i|x_i]\\
        & = E[z_i\beta|x_i] + E[g(x_i)|x_i] + \underbrace{E[\epsilon_i|x_i]}_{(1)} = \boxtimes\\
        (1) E[\epsilon_i|x_i] & = E[E[\epsilon_i|x_i,z_i]|x_i] = E[0|x_i] = 0\\
        \boxtimes & = \beta E[z_i|x_i] + g(x_i)
    \end{align*}
    Now moving on to the second equation,
    \begin{align*}
        y_i - \underbrace{E[y_i|x_i]}_{\text{Nadaraya-Watson}} 
        & = z_i\beta + g(x_i) + \epsilon_i - E[z_i|x_i]\beta - g(x_i)\\
        & = [z_i - \underbrace{E[z_i|x_i]}_{\text{Nadaraya-Watson}}]\beta + \epsilon_i
    \end{align*}
    So we have,
    \begin{equation*}
        y_i - \widehat{E[y_i|x_i]} = [z_i - \widehat{E[z_i|x_i]}]\beta+\epsilon_i
    \end{equation*}
\end{proof}

If $E[y_i|x]$ and $E[z_i|x]$ are \Emph{known}, $\beta$ can be estimated by \KeyIdea{OLS}. 
However, if $E[y_i|x]$ and $E[z_i|x]$ are \Emph{replaced by kernel estimators} that converges sufficiently quickly,
their use will not affect the rate of convergence of the \KeyIdea{OLS} estimator of $\beta$
\begin{equation*}
    \hat{\beta} = \left[(Z-\widehat{E[Z|x]})'(Z-\widehat{E[Z|x]})\right]^{-1}(Z-\widehat{E[Z|x]})'(Y-\widehat{E[Z|x]})
\end{equation*}
This estimator is \Emph{consistent and asymptotically normal}.






% \begin{definition}
% Theorem
% \begin{lemma}
% \begin{Example}
% \begin{proof}
%\newcommand{\KeyIdea}  % 关键想法：浅绿底
%\newcommand{\Term}     % 专业名词：品牌深绿字
%\newcommand{\Emph}     % 强调词汇：红
%\newcommand{\Confuse}    % 不懂/存疑处
%\newcommand{\Info}     % 说明或注解